% ==========================================================================
%  Honeybee Integration - Build Session Transcript (thesis appendix)
%  Human-readable summary transcript of the working session in which an AI
%  agent integrated the BioDT honeybee app into ShinySwarm. Same style as
%  docs/shinyswarm-build-session.txt. Records conversation flow, decisions,
%  and outcomes -- not a byte-exact chat log. Reconstructed from the session
%  of 2026-06-12. \input or \include this as an appendix.
% ==========================================================================

\section{Appendix: Honeybee Integration Build Session}
\label{app:honeybee-session}

\textbf{Project:} Collaborative-RShiny (ShinySwarm), Master's thesis, UvA.\\
\textbf{Author:} Iva Gunjaca \quad \textbf{Agent:} Ona \quad \textbf{Date:} 2026-06-12.\\
\textbf{Topic:} Re-engineering the BioDT honeybee (Beekeeper pDT) use case into
the ShinySwarm REST and Kafka architectures, with a simulated compute backend.

This is a human-readable summary of the working session: the request flow, the
decisions taken, and the outcomes. The artifacts it produced live under
\texttt{Thesis-Project-Final-RESTAPI/shiny\_services/honeybee/},
\texttt{Thesis-Project-Final-Kafka/shiny\_services/honeybee/}, the two Spring
\texttt{DataInitializer.java} files, \texttt{CollaborationController.java}, the
\texttt{docker-compose.yml} / \texttt{nginx.conf} of both stacks, and
\texttt{docs/thesis/ai\_integration\_experiment.tex}.

% --------------------------------------------------------------------------
\subsection*{1. Request: assess and perform the integration}
% --------------------------------------------------------------------------

User asked the agent to read the thesis docs and notes and judge whether the
``BioDT beehive app'' could be integrated into ShinySwarm, noting that backend
values could be simulated in this environment. The agent read
\texttt{README.md}, the existing \texttt{biodt\_feasibility.tex}, the
architecture summary, and the \texttt{statesnap} sources, then inspected the
two architectures and the eight benchmark apps. It reported that the integration
was feasible and that the map and Monte~Carlo apps between them already covered
the honeybee interaction model (Leaflet location picking + stochastic
population simulation).

The user then set the scope: \emph{full proof of concept}, \emph{both}
architectures. A later clarification established the framing -- the experiment of
interest is ``how well can AI integrate an existing app,'' with this conversation
itself as the reference -- and that the write-ups should stay as lightweight
notes, not a polished chapter.

% --------------------------------------------------------------------------
\subsection*{2. Grounding in the host pattern}
% --------------------------------------------------------------------------

Before writing code the agent read two existing apps in full --
\texttt{shiny\_services/map/} and \texttt{shiny\_services/monte\_carlo/} (front
and back) -- plus \texttt{CollaborationController.java}. It extracted the
repeated ShinySwarm shape: front/back split; \texttt{identity()} parsed from the
URL query string; VIEWER/EDITOR role gating; a POST-then-poll relay (REST) or
Producer/Consumer (Kafka); and \texttt{appName}-based routing in the relay.

% --------------------------------------------------------------------------
\subsection*{3. Acquiring the real upstream source}
% --------------------------------------------------------------------------

Rather than reconstruct a lookalike UI from the feasibility notes, the agent
cloned \texttt{github.com/BioDT/biodt-shiny} and read the actual honeybee module
(\texttt{app/view/honeybee/*}, \texttt{app/logic/honeybee/*}). This corrected a
first-draft backend that had \emph{invented} parameter names: the real module
uses \texttt{N\_INITIAL\_BEES}, \texttt{N\_INITIAL\_MITES\_HEALTHY},
\texttt{N\_INITIAL\_MITES\_INFECTED}, \texttt{HoneyHarvesting},
\texttt{VarroaTreatment}, \texttt{DroneBroodRemoval}, plus a
year/\texttt{DaysLimit} window and an editable habitat lookup table.

% --------------------------------------------------------------------------
\subsection*{4. Pinning the I/O contract}
% --------------------------------------------------------------------------

Inputs were traced to \texttt{beekeeper\_param.R} and
\texttt{beekeeper\_runsimulation.R}; outputs to
\texttt{honeybee\_beekeeper\_plot.R}, whose result table uses the exact columns
\texttt{date}, \texttt{weather}, \texttt{TotalIHbees + TotalForagers}, and
\texttt{(honeyEnergyStore / ( ENERGY\_HONEY\_per\_g * 1000 ))}. The simulated
backend was written to consume those inputs and emit those columns verbatim.

% --------------------------------------------------------------------------
\subsection*{5. Implementation (both architectures)}
% --------------------------------------------------------------------------

\begin{itemize}
  \item \textbf{REST:} \texttt{shiny\_back\_beehave.r} (Plumber, simulated
    BEEHAVE) and \texttt{shiny\_front\_beehave.r} (BioDT UI flattened from Rhino
    \texttt{box} modules to plain Shiny, wired to the relay with 500\,ms polling).
  \item \textbf{Kafka:} the same UI as a Producer/Consumer app, and a consumer
    worker that runs the same model and produces results to the \texttt{output}
    topic -- the better fit for BEEHAVE's long-running jobs.
  \item \textbf{Collaboration:} apiary coordinates, lookup-table edits, and run
    results shared across users; VIEWER/EDITOR gating reused from the benchmark
    apps.
  \item \textbf{statesnap:} reproducible checkpoints capturing the
    \emph{computed} (non-deterministic) result, so a restore reproduces the exact
    trajectory that input-only bookmarking cannot.
\end{itemize}

% --------------------------------------------------------------------------
\subsection*{6. Wiring and invariants}
% --------------------------------------------------------------------------

The new app was registered as App~\#9 in both \texttt{DataInitializer.java}
files -- appended, not reordered, because IDs~1--8 are used by the k6 benchmarks.
REST routing added \texttt{"Honeybee" -> http://shiny-back-beehave:8000/state}
to \texttt{APP\_ROUTES}; Kafka routing is by \texttt{appName} in the message.
Service entries, an nginx server block, and proxy ports were added to both
stacks (REST \texttt{:8089}, Kafka \texttt{:11088}). An attribution/provenance
header was added to each new file.

% --------------------------------------------------------------------------
\subsection*{7. Verification}
% --------------------------------------------------------------------------

R~4.3.3 was installed. All four R artifacts and the \texttt{statesnap} sources
\texttt{parse()} cleanly. The simulation function was executed: it returns a
$366\times4$ table for a 365-day run with BioDT's column names, and is
genuinely non-deterministic (peak colony size $18{,}995$ vs $19{,}323$ across two
seeds) -- the property that makes the \texttt{statesnap} reproducibility argument
meaningful. The full Docker stack was \emph{not} brought up; the
Java/compose/nginx wiring is verified by inspection, not by a live end-to-end
run.

% --------------------------------------------------------------------------
\subsection*{8. Attribution disclaimer added}
% --------------------------------------------------------------------------

At the user's request, an attribution/disclaimer was added to all honeybee
artifacts, the experiment notes, and the commit message: the honeybee
application originates with the \textbf{LTER LIFE} project and is used here with
their permission for this thesis; all credit for the original application and its
scientific model belongs to the LTER LIFE / BioDT authors.

% --------------------------------------------------------------------------
\subsection*{9. Code review (independent pass)}
% --------------------------------------------------------------------------

A review pass found, most importantly, a \textbf{behavioural bug the builder
introduced}: both backends only handled the \texttt{RUN\_SIMULATION} command and
silently dropped the \texttt{COORDS} and \texttt{LOOKUP} state-sharing messages,
so two of the three advertised collaboration features did not propagate
end-to-end (and the REST variant could overwrite a stored result in Redis). The
reference \texttt{map} app avoids this by echoing inbound events back with a
timestamp. The review also flagged \texttt{allow\_rds = TRUE} as unsafe for the
cross-user restore scenario, the \texttt{Map.of} ten-pair ceiling, and a
possible missing \texttt{leaflet.extras}/\texttt{DT} dependency in the worker
image. These were recorded as findings; the bug fix was left pending at the
user's discretion. Notably, the defect surfaced only under review -- evidence
both for the architecture's regularity (the agent reproduced the contract and
patterns in one session) and for the limits of unreviewed AI integration.

% --------------------------------------------------------------------------
\subsection*{10. Final state of deliverables}
% --------------------------------------------------------------------------

\begin{itemize}
  \item \texttt{Thesis-Project-Final-RESTAPI/shiny\_services/honeybee/shiny\_\{front,back\}\_beehave.r}
  \item \texttt{Thesis-Project-Final-Kafka/shiny\_services/honeybee/shiny\_\{front,back\}\_beehave.r}
  \item \texttt{CollaborationController.java} (REST route), both
    \texttt{DataInitializer.java} (App~\#9)
  \item \texttt{docker-compose.yml} + \texttt{nginx.conf} (both stacks)
  \item \texttt{docs/thesis/ai\_integration\_experiment.tex} (experiment notes)
  \item \texttt{docs/thesis/ai\_integration\_prompt.md} (reusable prompt)
\end{itemize}

\noindent\textbf{Verification:} 4/4 R files parse; simulation runs and is
non-deterministic with BioDT's output schema; full stack not run end-to-end.\\
\textbf{Git:} committed on branch \texttt{honeybee-integration} and pushed;
nothing merged.
